
\subsection{The boundary-field chain and the additional estimate}
\label{obc:setup}

We now extend the diagrammatic doubling argument from periodic rings
to open chains with endpoint magnetic fields. Consider the spin-1
Heisenberg chain with equal endpoint fields, \fieldsame,
\begin{equation}
 H_n^{\same}
 =\sum_{i=1}^{n-1}\bm S_i\cdot\bm S_{i+1}
 -h(S_1^z+S_n^z),\qquad h=\frac35,
 \label{obc:model}
\end{equation}
where $\sum_\alpha(S^\alpha)^2=2I$ and the exchange coupling is one.
We give a diagrammatic proof of the following theorem from
Ref.~\cite{OpenAI2026Boundaryfield}, using its certified finite inputs.

\begin{theorem}[Boundary-field gap]
\label{obc:theorem}
For every odd $n\ge1921$, $H_n^{\same}$ has a unique ground state and
\[
 \Delta_n>\frac{\log 10}{392}\simeq0.005874.
\]
\end{theorem}

Without the endpoint fields, the effective spin-$1/2$ modes at the
two ends produce four low-lying states whose splitting is exponentially
small in the chain length~\cite{Kennedy1990Exact}. Applying equal fields on an odd chain lifts this
boundary degeneracy~\cite{OpenAI2026Boundaryfield}.

Our strategy is to bound the thermal purity of the open chain at a
scale $\beta$ proportional to the length, just as in the periodic
argument. List the energies of $H_n^{\same}$ with multiplicity and let $w_0\ge w_1\ge\cdots$
be their Gibbs weights. With
$\Pt^{\same}(n,\beta)=\sum_m w_m^2$ and
$\Delta_n=E_1(n)-E_0(n)$, we have
\begin{equation}
 w_0\ge\Pt^{\same}(n,\beta),\qquad
 e^{-\beta\Delta_n}
 \le\frac{1-\Pt^{\same}(n,\beta)}{\Pt^{\same}(n,\beta)}.
 \label{obc:target}
\end{equation}
A thermal purity above $1/2$ forces a unique ground state and bounds the
gap from below. What changes in the open-chain setting is how we
propagate the thermal purity across doubled lengths and inverse
temperatures.

To see where the new difficulty enters, consider the spatial transfer
operator at a fixed temperature. Suppose that $X_\beta$ has a simple
positive leading eigenvalue $\lambda$, \specbars[mpored]{.9,.06,.04}, and
write the normalized transfer operator as
\begin{equation}
 T:=X_\beta/\lambda=\Pi+R,\qquad
 \Pi=\ket{\psi}\bra{\psi},\qquad
 \Pi R=R\Pi=0,\quad \|\psi\|=1.
 \label{obc:split}
\end{equation}
In a periodic ring, the trace weights the leading projector $\Pi$ by one.
In an open chain, the chain ends are contracted with boundary vectors,
so the leading contribution is weighted instead by how much of $v$ lies
along $\psi$, \overlapglyph:
\begin{equation}
 \underbrace{\tnchain[ring]{\Pi}}_{\trv\Pi=1},
 \qquad
 \underbrace{\tnchain[capl=\bar v,capr=v]{\Pi}}_
 {c=|\langle\psi,v\rangle|^2}.
 \label{obc:leading-diagram}
\end{equation}
Because $c$ is not fixed by normalization and could a priori vanish,
squaring in space no longer closes on its own. For an even length $m$,
the normalized boundary matrix element splits into the leading overlap
and a nonnegative remainder,
\begin{align}
 B_m
 &=\tnchain[capl=\bar v,capr=v]{T^m}\notag \\
 &=\tnchain[capl=\bar v,capr=v]{\Pi}
   +\tnchain[capl=\bar v,capr=v]{R^m}
 =c+D_m,\notag \\
 D_m&=\|R^{m/2}v\|^2\ge0.
 \label{obc:boundary-split}
\end{align}
To show that the leading eigenvector dominates $B_m$, we must bound
the remainder's share of $B_m$,
\begin{equation}
 e_m:=\frac{D_m}{B_m}=1-\frac{c}{B_m},
 \label{obc:defect}
\end{equation}
the grey fraction of \stackbar{.14} at the seed value $e_{142}\le.14$
of Eq.~\eqref{obc:seeds}, rather than the absolute remainder $D_m$
alone. Increasing the chain length suppresses $D_m$ by powers of
$\|R\|$ while leaving $c$ unchanged.

When we double the inverse temperature, however, both the boundary vector $v$ and the leading eigenvector $\psi$ change. The overlap $c$ at the old temperature does not by itself bound the new overlap $c'$. The key idea is to compare two even lengths at the new temperature. Writing primes for quantities at $2\beta$, the difference
\[
 B_k'-B_m'=D_k'-D_m'
\]
cancels the unknown overlap $c'$ because $c'$ is independent of length.
An exact thermal-purity bridge then relates the ratio $B_k'/B_m'$ back to
known quantities at the old temperature.

The proof therefore proceeds in four stages:
\begin{enumerate}
 \item In Sec.~\ref{obc:caps}, we express the open-chain partition functions as boundary contractions of the continuous-time transfer operator. Putting the same boundary vector at both ends produces an auxiliary opposite-field chain with nonnegative remainders at even lengths, while a $\pi$ spin rotation at one end restores equal fields.
 \item In Sec.~\ref{obc:ring}, we use the periodic doubling recurrence to control the bulk transfer spectrum and spectral ratio across all scales, and we state the four finite seed bounds needed to start the open-chain recurrence.
 \item In Secs.~\ref{obc:steps} and~\ref{obc:initialization}, we combine the two-length boundary comparison, length doubling, and thermal squaring into a coupled recurrence for $e_m$ and the even-length opposite-field thermal purity, and we run three explicit initialization steps before entering the small-defect induction.
 \item In Sec.~\ref{obc:comparison}, we pass from even opposite-field chains to odd equal-field chains by bounding the signed remainder in absolute value, completing the proof of Theorem~\ref{obc:theorem}.
\end{enumerate}

\subsection{Boundary vectors and boundary conditions}
\label{obc:caps}

We first represent the open-chain partition functions as transfer-operator
matrix elements. We use the continuous-time bond expansion of
Ref.~\cite{OpenAI2026Boundaryfield}: we sum over the times of bond
insertions instead of placing them on a finite Trotter grid.
In the Cartesian basis, $G_\alpha=iS^\alpha$ is real antisymmetric and
\[
 -\bm S_i\cdot\bm S_{i+1}
 =\sum_{\alpha=x,y,z}G_{\alpha,i}G_{\alpha,i+1}.
\]
A bond history $a$ is an ordered list of events $(\tau,\alpha)$ with
$0<\tau<\beta$. The measure $\mu_\beta$ sums the labels and integrates
the ordered times; its total mass is $e^{3\beta}$, including unit mass
on the empty history. The local tensor is
\begin{align}
 W_\beta^{ss'}(a,b)
 &=e^{-\kappa\beta}
 \bra{s}\mathcal T_\tau
 \prod_{(\tau,\alpha)\in a\cup b}G_\alpha\ket{s'},\notag \\
 \kappa&=\frac{700741}{500000}.
 \label{obc:local}
\end{align}
Time ordering $\mathcal T_\tau$ places later events to the left
according to the interleaved times in $a\cup b$, and tracing the
physical indices defines the real symmetric transfer kernel
$X_\beta(a,b)$ on $L^2(\mu_\beta)$.

The energy shift $\kappa$ is the one used in
Ref.~\cite{OpenAI2026Boundaryfield}. Throughout this section, $Z_n$ and all
open-chain partition functions include the shift $\kappa nI$;
it places the leading transfer eigenvalue near one at the seed scales
and cancels from every purity and spectral gap.

To represent a boundary field along $z$, we cap the virtual bond history
with the boundary vector
\begin{equation}
 v_\beta(a)=
 \begin{cases}
 (ih)^{|a|},&\text{all event labels are }z,\\
 0,&\text{otherwise},
 \end{cases}
 \qquad \|v_\beta\|^2=e^{\beta h^2}.
 \label{obc:vector}
\end{equation}
Its value on the empty history is one, and its norm follows by summing
$(\beta h^2)^k/k!$ over the number $k$ of $z$ events.

The cap inserts boundary events into the same time-ordered product
as the bond events. Summing their counts and integrating their times
adds the corresponding field term to the Hamiltonian.
We now check the signs of the two endpoint fields produced by $v_\beta$.
The conjugated left cap supplies $(-ih)G_z=+hS^z$ in the exponential,
which corresponds to the Hamiltonian term $-hS_1^z$. The unconjugated right cap supplies
$(ih)G_z=-hS^z$, corresponding to $+hS_n^z$. Contracting with the same boundary vector $v_\beta$ at both
ends therefore gives the auxiliary opposite-field chain \fieldopp,
\begin{equation}
 H_n^{\opp}=H_n^{\same}+2hS_n^z.
 \label{obc:opposite}
\end{equation}
To obtain equal endpoint fields, let $P=\operatorname{diag}(1,-1,-1)$ be the spin rotation by $\pi$
about the $x$ axis. Its virtual action $U_P$ negates each $z$ event, which reverses
the right endpoint field and produces $H_n^{\same}$, \fieldsame. In diagram form,
\begin{align}
 e^{-\beta(H_n^{\opp}+\kappa nI)}&=\obcW,\notag \\
 Y_n(\beta)
 :=\Tr e^{-\beta(H_n^{\opp}+\kappa nI)}
 &=\obccap{X_\beta^n}{v_\beta},\notag \\
 Y_n^{\same}(\beta)
 :=\Tr e^{-\beta(H_n^{\same}+\kappa nI)}
 &=\obccap{X_\beta^n}{U_Pv_\beta}.
 \label{obc:capping}
\end{align}
For an $n$-site chain, site $i$ carries $W_\beta(a_{i-1},a_i)$, with
internal bond histories $a_1,\ldots,a_{n-1}$ and boundary histories
$a_0,a_n$ contracted with $\bar v_\beta,v_\beta$. Taking the physical
trace at each of the $n$ sites produces $X_\beta^n$ and gives
$Y_n(\beta)$ in the second line; inserting
$U_P$ at the right cap gives the equal-field partition function $Y_n^{\same}(\beta)$ in the third.

Spin rotations act on bond histories by real orthogonal operators that commute
with $X_\beta$, so $[U_P,X_\beta]=0$. Regrouping the expansion into these site-by-site
contractions is exact: before the energy shift, the absolute traced series
is bounded by
$3^n e^{\beta[3(n-1)+2h]}<\infty$, so Fubini's theorem justifies
contracting the time-ordered expansion site by site.

Having expressed both open chains through $X_\beta$ and $v_\beta$, we work first with the opposite-field partition function $Y_m(\beta)$ at even lengths $m$, where
Eq.~\eqref{obc:boundary-split} decomposes the boundary contraction into a sum of nonnegative terms.
We will return to the odd equal-field chain $Y_n^{\same}(\beta)$ in Sec.~\ref{obc:comparison}, using an absolute remainder
bound that does not require odd powers of the transfer operator to be positive.

\subsection{Finite inputs and ring estimates}
\label{obc:ring}

We now identify the finite inputs required to drive the open-chain recurrence across the dyadic scales
\begin{equation}
 \beta_j=49\,2^{j-2},\qquad n_j=120\,2^{j-2},\qquad
 M_j=2n_j,\qquad j\ge2,
 \label{obc:scales}
\end{equation}
where $n_j$ is a ring length and $M_j$ is an even open-chain length.
We retain the index $j=2$ for the certified starting point of
Ref.~\cite{OpenAI2026Boundaryfield}, at $\beta_2=49$.
At each scale $j$, the two-length boundary comparison and the length-doubling step require three bulk properties of the transfer operator $X_{\beta_j}$: a simple positive leading eigenvalue $\lambda_j$ invariant under $U_P$, a bound on the subleading norm $\|R_j\|^{n_j}$ to suppress the boundary remainder across the length difference $M_j-k_j$, and a lower bound on the bulk normalization ratio $r_j=\lambda_{j+1}/\lambda_j^2$ relating consecutive temperatures. All three bulk properties follow at every $j\ge2$ from the periodic spatial and thermal purities $\Ps(n_2,\beta_2)$ and $\Pt(M_2,\beta_2)$ at the initial scale $j=2$ ($\beta_2=49$, $n_2=120$, $M_2=240$). To start the open-chain recurrence itself, we also need an initial bound on the boundary defect $e_{k_2}(2)$ at a shorter even length $k_2=142<M_2$, together with a lower bound on the opposite-field thermal purity $\Pt^{\opp}(M_2,\beta_2)=Y_{M_2}(2\beta_2)/Y_{M_2}(\beta_2)^2$.

Ref.~\cite{OpenAI2026Boundaryfield} certifies these four finite seed bounds at $j=2$:
\begin{align}
 \Ps(120,49)&\ge.948,\qquad \Pt(240,49)\ge.948,\notag \\
 e_{142}(2)&\le.14,\qquad \Pt^{\opp}(240,49)\ge.79.
 \label{obc:seeds}
\end{align}
Here $\Ps(n,\beta)=Z_{2n}(\beta)/Z_n(\beta)^2$ and
$\Pt(n,\beta)=Z_n(2\beta)/Z_n(\beta)^2$ concern the periodic ring, while
$e_m(j)$, defined by Eq.~\eqref{obc:defect} at $\beta_j$, and
$\Pt^{\opp}(m,\beta)=Y_m(2\beta)/Y_m(\beta)^2$ concern the opposite-field open chain.

We first use the two ring inputs in Eq.~\eqref{obc:seeds} to control the bulk transfer operator at every scale $j\ge2$. By the periodic doubling argument, the ring purities satisfy
\begin{align}
 &\Ps(n_j,\beta_j),\ \Pt(M_j,\beta_j)\ge1-\varepsilon_j,\notag \\
 \varepsilon_2&=.052,\qquad
 \varepsilon_{j+1}=6\varepsilon_j^2.
 \label{obc:ring-purities}
\end{align}
To verify the factor of $6$ in this scalar recurrence, recall the squaring inequality:
if a positive trace-one operator of purity $P$ is squared and
renormalized, which concentrates its weights,
$\specbars{.5,.3,.2}\mapsto\specbars{.66,.24,.1}$, its new purity $P'$
satisfies
\begin{equation}
 1-P'\le f(1-P),\qquad
 f(z)=\frac{z^2}{2(1-z)^2}.
 \label{obc:square}
\end{equation}
The periodic coupling identity first gives
$\Ps(n,2\beta)\ge(1-\varepsilon)^3$. For
$0\le\varepsilon\le.052$, we have
\begin{align}
 f(1-(1-\varepsilon)^3)
 &=\frac{\varepsilon^2}{2}
   \left[\sum_{s=1}^3(1-\varepsilon)^{-s}\right]^2\notag \\
 &\le6\varepsilon^2\le\varepsilon.
 \label{obc:six}
\end{align}
After spatial squaring bounds the new spatial defect by $6\varepsilon^2$, the second periodic coupling identity gives
$\Pt(4n,\beta)\ge(1-\varepsilon)^2(1-6\varepsilon^2)
\ge(1-\varepsilon)^3$. Thermal squaring then bounds the new thermal defect by the same
$6\varepsilon^2$, proving
Eq.~\eqref{obc:ring-purities} by induction on $j\ge2$ without waiting for the smaller open-chain threshold established below.

From the ring purities in Eq.~\eqref{obc:ring-purities}, we now extract the bulk spectral properties of $X_{\beta_j}$ needed for the open chain. Write $\lambda_j$ for the leading eigenvalue of $X_{\beta_j}$, and decompose
$T_j=X_{\beta_j}/\lambda_j=\Pi_j+R_j$ as in Eq.~\eqref{obc:split}. We claim that for every $j\ge2$,
\begin{align}
 \lambda_j>0\text{ is simple},\qquad
 \|R_j\|^{n_j}&\le d(\varepsilon_j),\notag \\
 d(\varepsilon)&=\frac{\varepsilon}{2-3\varepsilon},
 \qquad \|R_j\|<1,\notag \\
 r_j:&=\frac{\lambda_{j+1}}{\lambda_j^2},\qquad
 0<r_j\le1,\notag \\
 r_j^{M_j}&\ge(1-\varepsilon_j)(1-\varepsilon_{j+1}),
 \qquad U_P\Pi_j=\Pi_j.
 \label{obc:ring-bounds}
\end{align}

We verify each assertion in Eq.~\eqref{obc:ring-bounds} in turn. For any probability distribution
of purity at least $1-\varepsilon>1/2$, let $p_0$ be its largest weight
and put $z=1-p_0$. Since $p_0$ bounds the purity from above, $z<1/2$, and $1-\varepsilon\le p_0^2+(1-p_0)^2=1-2z(1-z)$ gives $2z(1-z)\le\varepsilon$. Hence
\[
 z\le\varepsilon,\qquad
 z\le\frac{\varepsilon}{2(1-\varepsilon)},\qquad
 \frac{z}{1-z}\le d(\varepsilon).
\]
Applying this estimate to the normalized weights
$p_i=|\lambda_i(\beta_j)|^{n_j}/Z_{n_j}(\beta_j)$ yields the subleading tail bound
\begin{equation}
 \sum_{i\ge1}
 \left|\frac{\lambda_i(\beta_j)}{\lambda_j}\right|^{n_j}
 \le d(\varepsilon_j).
 \label{obc:full-tail}
\end{equation}
Because $d(\varepsilon_j)<1$, the largest-modulus eigenvalue $\lambda_j$ is simple and $\|R_j\|^{n_j}\le d(\varepsilon_j)<1$. Moreover, $\lambda_j$ is positive: because $X_{\beta_j}$ is Hilbert--Schmidt and the remaining spectrum has strictly smaller spectral radius, a negative isolated leading
eigenvalue would make $\trv X_{\beta_j}^N$ negative for all sufficiently
large odd $N$, contradicting $Z_N(\beta_j)>0$.

Next, to see that $U_P\Pi_j=\Pi_j$, note that the three $\pi$ rotations about the coordinate axes and the cyclic
permutation of those axes commute with the real transfer kernel $X_{\beta_j}$.
Each acts by a sign on the real one-dimensional leading eigenspace. The order-three cyclic
permutation must act by $+1$, and it conjugates the three $\pi$ rotations into one
another. Their signs are therefore equal; since the product of two $\pi$
rotations is the third, the common sign is $+1$, which gives
$U_P\Pi_j=\Pi_j$.

Finally, to bound the bulk normalization ratio $r_j=\lambda_{j+1}/\lambda_j^2$, observe that $\Pt(N,\beta_j)\le1$ and taking even $N$th roots as $N\to\infty$
give $\lambda_{j+1}\le\lambda_j^2$, so $r_j\le1$. At length $M_j=n_{j+1}$, combining the leading
spatial weight at $\beta_{j+1}$ with the ring thermal purity at $\beta_j$ gives the lower bound
\begin{align}
 \lambda_{j+1}^{M_j}
 &\ge(1-\varepsilon_{j+1})Z_{M_j}(\beta_{j+1})\notag \\
 &\ge(1-\varepsilon_{j+1})(1-\varepsilon_j)
       Z_{M_j}(\beta_j)^2\notag \\
 &\ge(1-\varepsilon_{j+1})(1-\varepsilon_j)
       \lambda_j^{2M_j}.
 \label{obc:bulk}
\end{align}
Dividing by $\lambda_j^{2M_j}$ proves $r_j^{M_j}\ge(1-\varepsilon_j)(1-\varepsilon_{j+1})$.
With the bulk transfer estimates in Eq.~\eqref{obc:ring-bounds} established, we turn to the open-chain recurrence.

\subsection{Three steps for the open chain}
\label{obc:steps}

We now set up the open-chain quantities at scale $j$ and show how to advance them to scale $j+1$. Set $v_j=v_{\beta_j}$ and, for even $m\ge2$, define the normalized opposite-field partition function, its remainder, and the boundary defect by
\begin{align}
 c_j&=|\langle\psi_j,v_j\rangle|^2,\notag \\
 B_m(j)&=\frac{Y_m(\beta_j)}{\lambda_j^m}
       =c_j+D_m(j)>0,\notag \\
 D_m(j)&=\|R_j^{m/2}v_j\|^2,\qquad
 e_m(j)=\frac{D_m(j)}{B_m(j)}.
 \label{obc:normalized}
\end{align}
For even lengths $m\ge k$, absorbing $R_j^{k/2}$ into each boundary cap gives the spatial decay of the remainder,
\begin{align}
 D_m(j)
 &=\tnchain[capl={\bar y},capr=y]{R_j^{m-k}},
 \qquad y=R_j^{k/2}v_j,\notag \\
 0\le D_m(j)&\le\|R_j\|^{m-k}D_k(j).
 \label{obc:decay}
\end{align}
The remainder therefore decays geometrically with the extra length $m-k$,
\decayglyph[black!50], and in particular $D_m(j)$ is nonincreasing in
even $m$, so $c_j\le B_m(j)\le B_k(j)$.

To connect quantities at $\beta_j$ and $\beta_{j+1}=2\beta_j$, we use the opposite-field thermal purity. By the definition of $B_m(j)$ and $r_j=\lambda_{j+1}/\lambda_j^2$, we have the exact thermal-purity bridge
\begin{align}
 \Pt^{\opp}(m,\beta_j)
 &=\frac{Y_m(\beta_{j+1})}{Y_m(\beta_j)^2}\in(0,1],\notag \\
 B_m(j+1)
 &=\Pt^{\opp}(m,\beta_j)r_j^{-m}B_m(j)^2.
 \label{obc:bridge}
\end{align}
Multiplying the normalized open-chain Gibbs operators
$\rho_m^{\opp}(\beta)=e^{-\beta(H_m^{\opp}+\kappa mI)}/Y_m(\beta)$
in imaginary time gives
\begin{equation}
 \tncolumn{\rho_m^{\opp}(\beta_j),\rho_m^{\opp}(\beta_j)}
 =\Pt^{\opp}(m,\beta_j)
  \tncolumn{\rho_m^{\opp}(\beta_{j+1})}.
 \label{obc:thermal-diagram}
\end{equation}

Assuming bounds on a boundary defect $e_k(j)$ and on the thermal purity $\Pt^{\opp}(m,\beta_j)$ at scale $j$, we advance to scale $j+1$ in three steps:
first we bound the boundary defect $e_m(j+1)$ at the new temperature;
second we use that new boundary defect to double the chain length from $m$ to $2m$ in the thermal purity at $\beta_j$;
third we square in imaginary time to obtain $\Pt^{\opp}(2m,\beta_{j+1})$.

\subsubsection{Step 1: control the new boundary remainder}
\label{obc:move-one}

Our first task is to bound the boundary defect $e_m(j+1)$ at the new inverse temperature $\beta_{j+1}$ without knowing the new leading overlap $c_{j+1}$ in advance. Fix even lengths $m>k\ge2$, assume $e_k(j)<1$, and write primes for quantities at scale $j+1$. Choose a suppression factor $q<1$ satisfying
$\|R_{j+1}\|^{m-k}\le q$, so that Eq.~\eqref{obc:decay} at scale $j+1$ gives
$D_m'\le qD_k'$. Subtracting the normalized boundary contractions at lengths $k$ and $m$ at the same temperature $\beta_{j+1}$ cancels their common leading term $c'$:
\begin{align}
 &\tnchain[capl={\bar v'},capr={v'}]{(T')^k}
 -\tnchain[capl={\bar v'},capr={v'}]{(T')^m}\notag \\
 &\quad=
 \tnchain[capl={\bar v'},capr={v'}]{(R')^k}
 -\tnchain[capl={\bar v'},capr={v'}]{(R')^m}.
 \label{obc:cancellation-diagram}
\end{align}
That is, $B_k'-B_m'=D_k'-D_m'$. Combining $D_k'=D_m'+B_k'-B_m'$ with $D_m'\le qD_k'$ gives
\begin{align}
 (1-q)D_m'&\le q(B_k'-B_m'),\notag \\
 e_m'&\le\frac{q}{1-q}\left(\frac{B_k'}{B_m'}-1\right).
 \label{obc:new-defect}
\end{align}
The unknown overlap $c'$ has dropped out of the difference $B_k'-B_m'$.

We can now bound the new-temperature ratio $B_k'/B_m'$ in terms of scale-$j$ quantities by applying the bridge identity~\eqref{obc:bridge} at both lengths $k$ and $m$:
\begin{align}
 \frac{B_k'}{B_m'}
 &=\frac{\Pt^{\opp}(k,\beta_j)}{\Pt^{\opp}(m,\beta_j)}
   r_j^{m-k}\left[\frac{B_k(j)}{B_m(j)}\right]^2\notag \\
 &\le\frac{1}{\Pt^{\opp}(m,\beta_j)[1-e_k(j)]^2}.
 \label{obc:ratio}
\end{align}
Here we used $\Pt^{\opp}(k,\beta_j)\le1$, $r_j\le1$, and the lower bound
\[
 B_m(j)\ge c_j=[1-e_k(j)]B_k(j)>0
\]
at the old temperature. Equation~\eqref{obc:ratio} makes no assumption on the new overlap $c'$; once Eqs.~\eqref{obc:new-defect} and~\eqref{obc:ratio} yield
$e_m'<1$, it follows automatically that $c'=(1-e_m')B_m'>0$.

Notice why Eq.~\eqref{obc:new-defect} produces a quadratic improvement in the small-defect regime. It is a product of two small factors: the spatial decay factor $q/(1-q)$ over the length difference $m-k$, and the excess of the ratio $B_k'/B_m'$ above one.
Specifically, when $m=M_j$, $k=M_j/2$, and $\varepsilon_j\le x\le.01$,
the ring bound~\eqref{obc:ring-bounds} at scale $j+1$ gives
\[
 \|R_{j+1}\|^{M_j}\le d(\varepsilon_{j+1})
 \le d(6x^2).
\]
Across the half-chain $m-k=M_j/2$, we may therefore take
\begin{equation}
 q=\sqrt{d(\varepsilon_{j+1})}
 \le\sqrt{\frac{6x^2}{2-18x^2}}\le2x.
 \label{obc:half-chain}
\end{equation}
Separately, when $e_k(j)\le x$ and
$\Pt^{\opp}(m,\beta_j)\ge1-x$, Eq.~\eqref{obc:ratio} bounds the excess above one by
\begin{equation}
 \frac{B_k'}{B_m'}-1\le(1-x)^{-3}-1=O(x).
 \label{obc:ratio-excess}
\end{equation}
Multiplying Eqs.~\eqref{obc:half-chain} and~\eqref{obc:ratio-excess} bounds $e_m'$ by $O(x^2)$. The subtraction of one in Eq.~\eqref{obc:new-defect}, coming from the cancellation of $c'$ in Eq.~\eqref{obc:cancellation-diagram}, supplies this second factor of $x$.

\subsubsection{Step 2: double the length}
\label{obc:move-two}

Having bounded the boundary defect $e_M(j+1)$ at the new temperature, we next double the open-chain length from $M$ to $2M$ in the thermal purity at the old temperature $\beta_j$.
For any even $M$, dividing the bridge identity~\eqref{obc:bridge} at length $2M$ by that at length $M$ gives
\begin{align}
 \frac{\Pt^{\opp}(2M,\beta_j)}{\Pt^{\opp}(M,\beta_j)}
 &=r_j^M\frac{B_{2M}(j+1)}{B_M(j+1)}
   \left[\frac{B_M(j)}{B_{2M}(j)}\right]^2\notag \\
 &\ge r_j^M[1-e_M(j+1)].
 \label{obc:length}
\end{align}
In the last factor, monotonicity at the old temperature gives $B_M(j)\ge B_{2M}(j)$, so $[B_M(j)/B_{2M}(j)]^2\ge1$.
In the middle factor, however, $B_{2M}(j+1)\le B_M(j+1)$, so we must use the new-temperature boundary defect bound from Step~1:
\[
 B_{2M}(j+1)\ge c_{j+1}
 =[1-e_M(j+1)]B_M(j+1).
\]
Equation~\eqref{obc:length} therefore retains both the bulk spectral-ratio factor $r_j^M$, controlled by Eq.~\eqref{obc:ring-bounds}, and the new-temperature boundary factor
$1-e_M(j+1)$.

\subsubsection{Step 3: square in imaginary time}
\label{obc:move-three}

Once Step~2 bounds $\Pt^{\opp}(2M,\beta_j)$ at the doubled length $2M$, we double the inverse temperature from $\beta_j$ to $\beta_{j+1}$ by squaring the open-chain Gibbs state in imaginary time.
Just as in the periodic argument, Eqs.~\eqref{obc:thermal-diagram} and~\eqref{obc:square} give
\begin{equation}
 1-\Pt^{\opp}(2M,\beta_{j+1})
 \le f(1-\Pt^{\opp}(2M,\beta_j)).
 \label{obc:thermal-square}
\end{equation}
Notice why the three steps must be performed in this order: because $\Pt^{\opp}(M,\beta_j)$ already connects $\beta_j$ and $\beta_{j+1}$, doubling the length in Step~2 requires the new-temperature boundary defect $e_M(j+1)$ from Step~1 before we can square in imaginary time in Step~3.

\subsection{Initialization and the small-defect induction}
\label{obc:initialization}

We now assemble the three steps of Sec.~\ref{obc:steps} into a
recurrence across $j\ge2$. The initial boundary estimate in
Eq.~\eqref{obc:seeds} is certified at $142$ sites. We use that length
for the first comparison and the preceding target length thereafter:
\begin{equation}
 k_2=142,\qquad k_{j+1}=M_j.
 \label{obc:k}
\end{equation}
Since $M_{j+1}=2M_j$, this choice gives $k_j=M_j/2$ for all $j\ge3$. Let $E_j$ and $\eta_j$
denote upper bounds on $e_{k_j}(j)$ and $1-\Pt^{\opp}(M_j,\beta_j)$, respectively.
By Eqs.~\eqref{obc:ring-bounds}, \eqref{obc:new-defect}, \eqref{obc:ratio}, \eqref{obc:length}, and~\eqref{obc:thermal-square}, the updates from scale $j$ to scale $j+1$ take the form
\begin{align}
 F(q,e,t)&=\frac{q}{1-q}
             \left[\frac{1}{(1-e)^2t}-1\right],\notag \\
 q_j&=d(\varepsilon_{j+1})^{(M_j-k_j)/M_j},\notag \\
 E_{j+1}&\ge F(q_j,E_j,1-\eta_j),\notag \\
 A_j&=(1-\eta_j)(1-\varepsilon_j)
       (1-\varepsilon_{j+1})(1-E_{j+1}),\notag \\
 \eta_{j+1}&\ge f(1-A_j),
 \label{obc:updates}
\end{align}
valid whenever $E_j,\eta_j,E_{j+1},\eta_{j+1}\in[0,1)$.
The product $A_j$ bounds $\Pt^{\opp}(2M_j,\beta_j)$ from below by
Eq.~\eqref{obc:length}; thermal squaring then gives the update for
$\eta_{j+1}$.
Here the exponent $(M_j-k_j)/M_j$ in $q_j$ comes from $n_{j+1}=M_j$ and the new-temperature
ring remainder bound $\|R_{j+1}\|^{M_j}\le d(\varepsilon_{j+1})$. Because $F(q,e,t)$ increases with $q$ and $e$ and decreases with $t$,
any upper bound on $q_j$ may be substituted into the update for $E_{j+1}$.

At the initial scale $j=2$, the seed defects $E_2=.14$ and $\eta_2=.21$ from Eq.~\eqref{obc:seeds} are too large to enter the analytical small-defect induction ($x\le.01$) directly. We therefore carry out three explicit rational steps from $j=2$ to $j=5$. Starting with
$\varepsilon_2=.052$, $E_2=.14$, and $\eta_2=.21$, Eq.~\eqref{obc:updates} produces the outward-rounded rational bounds in the following table:
\begin{center}
\begin{tabular}{@{}rrrr@{}}
\toprule
$j$ & $E_j$ & $\eta_j$ & upper bound for $q_j$\\
\midrule
2 & .14 & .21 & .142\\
3 & .12 & .155 & .0282\\
4 & .016 & .026 & .0028\\
5 & .001 & .001 & ---\\
\bottomrule
\end{tabular}
\end{center}
All decimals in this table denote exact rationals. Both defects collapse
along the three steps, $E_j$ as \specbars[black!45]{1,.86,.11,.007} and
$\eta_j$ as \specbars{1,.74,.12,.005} ($j=2,\dots,5$, each relative to
its value at $j=2$): the first step gains little because $q_2\le.142$ is
still large, whereas once $q_j$ is small the defects square. At $j=2$,
the first suppression bound uses $k_2=142$ and $M_2=240$, so that
\[
 \frac{M_2-k_2}{M_2}=\frac{49}{120},\qquad
 d(\varepsilon_3)^{49}<.142^{120}.
\]
At $j=3$ and $j=4$, we have $(M_j-k_j)/M_j=1/2$, so we verify the upper bounds on $q_j$ by squaring $.0282$ and $.0028$ and comparing with
$d(\varepsilon_4)$ and $d(\varepsilon_5)$.
Substituting these values into Eq.~\eqref{obc:updates} verifies each row of the table by exact rational arithmetic. Together with
$\varepsilon_5<.000015$, the $j=5$ row brings all three defects below $.01$, establishing the base case of the small-defect induction:
\begin{align}
 \varepsilon_j,\ e_{k_j}(j),\
 1-\Pt^{\opp}(M_j,\beta_j)&\le x_j,\qquad j\ge5,\notag \\
 x_5&=.01,\qquad x_{j+1}=10x_j^2,\qquad
 x_j=10^{-1-2^{j-5}}.
 \label{obc:induction}
\end{align}

To prove the induction step in Eq.~\eqref{obc:induction}, assume that
the three bounds hold at scale $j\ge5$ and write $x=x_j\le.01$. One
step from scale $j$ to scale $j+1=j'$ ($M_{j+1}=2M_j$) proceeds along
two coupled lanes:
\begin{proofdiagram}
\begin{prooflanes}
 \lanebound{b}{e0}{0}{0}{$e_{M_j/2}(j)\le x$}
 \lanebound{t}{t0}{1}{0}{$\Pt^{\opp}(M_j,\beta_j)\ge1-x$}
 \lanebound{b}{e1}{0}{1}{$e_{M_j}(j+1)\le10x^2$}
 \lanebound{to}{t1}{1}{2}{$\Pt^{\opp}(2M_j,\beta_j)\ge1-2x-16x^2$}
 \lanebound{t}{t2}{1}{3}{$\Pt^{\opp}(2M_j,\beta_{j+1})\ge1-10x^2$}
 \lanemove{e0}{e1}{1}{\eqref{obc:new-defect}}
 \laneuse{t0}{1}
 \lanemove{t0}{t1}{2}{\eqref{obc:length}}
 \laneuse{e1}{2}
 \lanemove[mpoblue]{t1}{t2}{3}{square}
\end{prooflanes}
\end{proofdiagram}
For Step~1, Eqs.~\eqref{obc:half-chain} and~\eqref{obc:ratio-excess}
give $q_j\le2x$ and
\begin{equation}
 F(2x,x,1-x)
 \le\frac{6x^2}{(1-2x)(1-x)^3}<7x^2,
 \label{obc:boundary-step}
\end{equation}
so $E_{j+1}=10x^2$ bounds $e_{M_j}(j+1)$. For Step~2, substituting
$\eta_j\le x$, $\varepsilon_j\le x$, $\varepsilon_{j+1}\le6x^2$, and
$E_{j+1}\le10x^2$ into $A_j$ gives $A_j\ge1-2x-16x^2$. For Step~3,
since $0<x\le.01$, thermal squaring yields
\begin{equation}
 f(2x+16x^2)<3x^2,
 \label{obc:thermal-step}
\end{equation}
so $\eta_{j+1}=10x^2$ bounds $1-\Pt^{\opp}(2M_j,\beta_{j+1})$, and
$\varepsilon_{j+1}\le6x^2\le10x^2$ closes the induction at scale $j+1$.

\subsection{Equal fields, odd lengths, and the gap}
\label{obc:comparison}

We have established the small-defect bounds~\eqref{obc:induction} for even-length opposite-field chains at every scale $j\ge5$.
It remains to transfer these bounds to the target equal-field Hamiltonian $H_n^{\same}$ at arbitrary odd lengths $n\ge M_5+1=1921$ and extract the spectral gap.

Define the normalized equal-field partition function by
\[
 Q_j(n)=\frac{Y_n^{\same}(\beta_j)}{\lambda_j^n}.
\]
Because $U_P$ commutes with $X_{\beta_j}$ and satisfies $U_P\Pi_j=\Pi_j$ by Eq.~\eqref{obc:ring-bounds}, inserting $U_P$ at the right cap preserves the leading contribution $c_j$:
\begin{equation}
 \tnchain[last=sym]{\Pi_j,U_P}=\tnchain{\Pi_j}.
 \label{obc:twist}
\end{equation}
For any even $M\le n$, we split off $R_j^{M/2}$ on each side of the remainder contraction to obtain
\begin{align}
 Q_j(n)-c_j
 &=\tnchain[capl={\bar v_j},capr={v_j}]
 {R_j^{M/2},R_j^{n-M}U_P,R_j^{M/2}}\notag \\
 &=\langle y,R_j^{n-M}U_Py\rangle,\qquad
 y=R_j^{M/2}v_j.
 \label{obc:signed}
\end{align}
Because $\|R_j\|\le1$ and $U_P$ is orthogonal, we have $\|R_j^{n-M}U_P\|\le1$, and Cauchy--Schwarz bounds the signed remainder in absolute value by the even-length opposite-field remainder $D_M(j)$:
\begin{equation}
 |Q_j(n)-c_j|\le\|y\|^2=D_M(j).
 \label{obc:absolute}
\end{equation}
This inequality controls $Q_j(n)$ at odd $n$ without assuming that odd powers of $R_j$ are positive.

Now fix a scale $j\ge5$ and an odd length $n$ in the interval $M_j\le n\le2M_j$, and set
$M=M_j=k_{j+1}$. Applying Eq.~\eqref{obc:absolute} at $\beta_j$ and $\beta_{j+1}$ bounds $Q_j(n)$ from above and $Q_{j+1}(n)$ from below in terms of $B_M$:
\begin{align}
 Q_j(n)&\le c_j+D_M(j)=B_M(j),\notag \\
 Q_{j+1}(n)&\ge c_{j+1}-D_M(j+1)\notag \\
 &=[1-2e_M(j+1)]B_M(j+1)\notag \\
 &\ge(1-2x_{j+1})B_M(j+1)>0.
 \label{obc:comparison-bounds}
\end{align}
Combining Eq.~\eqref{obc:comparison-bounds} with the bridge identity~\eqref{obc:bridge} at length $M$, and using $r_j\le1$ and $n-M\le M$ together with the bulk ratio bound $r_j^M\ge(1-\varepsilon_j)(1-\varepsilon_{j+1})$ from Eq.~\eqref{obc:ring-bounds}, we obtain the equal-field thermal purity bound
\begin{align}
 \Pt^{\same}(n,\beta_j)
 &=r_j^n\frac{Q_{j+1}(n)}{Q_j(n)^2}\notag \\
 &\ge(1-2x_{j+1})\Pt^{\opp}(M,\beta_j)r_j^{n-M}\notag \\
 &\ge(1-2x_{j+1})(1-x_j)
       (1-\varepsilon_j)(1-\varepsilon_{j+1})\notag \\
 &\ge1-2x_j-26x_j^2\notag \\
 &\ge1-4x_j,
 \label{obc:final-purity}
\end{align}
where the coefficient $26=20+6$ collects
$2x_{j+1}=20x_j^2$ and $\varepsilon_{j+1}\le6x_j^2$.

\begin{proof}[Proof of Theorem~\ref{obc:theorem}]
The intervals $[M_j,2M_j]$ for $j\ge5$ cover every integer
$n\ge M_5=1920$. Given any odd $n\ge1921$, choose $j\ge5$ with $M_j\le n\le2M_j$.
Since $x_j\le x_5=.01$, Eq.~\eqref{obc:final-purity} gives
$w_0\ge\Pt^{\same}(n,\beta_j)\ge1-4x_j>1/2$, so $H_n^{\same}$ has a unique ground state. At this
length $n$ and finite inverse temperature $\beta_j$,
Eq.~\eqref{obc:target} gives
\begin{align}
 e^{-\beta_j\Delta_n}
 &\le\frac{4x_j}{1-4x_j}\notag \\
 &\le\frac{5}{12}\,10^{-2^{j-5}}
 <10^{-2^{j-5}},
 \label{obc:gap-final}
\end{align}
where we used $x_j=10^{-1}10^{-2^{j-5}}$ and
$1-4x_j\ge.96$, which yields the prefactor $.4/.96=5/12<1$.
Taking logarithms and dividing by $\beta_j=392\,2^{j-5}$ proves
$\Delta_n>(\log 10)/392$.
\end{proof}
